% Write the note inside the frame it belongs to:

\begin{frame}{Results}
  \begin{itemize}
    \item The gain holds across both attackers
  \end{itemize}
  \speakernote{Say the number out loud: 78 against 62.
    If asked about variance, point at the backup slide.}
\end{frame}

% Then choose a mode in the preamble. Two matter, and they are opposites:

\usetheme{lucid}                 % notes=off is the default:
                                 % notes are NOT compiled at all.
                                 % This is the file you hand out.

\usetheme[notes=pages]{lucid}    % each slide is followed by a typeset
                                 % notes page, in the deck's own colours.

% Two more, for presenting:
\usetheme[notes=only]{lucid}     % notes without slides -- a script
\usetheme[notes=second]{lucid}   % presenter view, notes on screen two

% Beamer's own \note{...} works too and obeys the same option.
% Tip: build twice -- once plain to present from, once with
% notes=pages to print.
